% ============================================================================
% LANGENTWURF - Sequenz 6a: Los lugares (Orte in der Stadt)
% ============================================================================
% TIER 1 (Vollständig) | Profil A (Gesamtschule heterogen)
% ============================================================================

\documentclass[11pt,a4paper]{article}

\usepackage[utf8]{inputenc}
\usepackage[T1]{fontenc}
\usepackage[ngerman]{babel}
\usepackage[left=2cm,right=2cm,top=2cm,bottom=2cm]{geometry}
\usepackage{helvet}
\usepackage{xcolor}
\usepackage{tabularx}
\usepackage{array}
\usepackage{enumitem}
\usepackage{fancyhdr}
\usepackage{lastpage}
\usepackage{tcolorbox}
\usepackage{booktabs}
\usepackage{amssymb}

\renewcommand{\familydefault}{\sfdefault}

\definecolor{spanisch}{HTML}{c41e3a}
\definecolor{grau}{HTML}{666666}
\definecolor{hellgrau}{HTML}{f0f0f0}
\definecolor{basis}{HTML}{4CAF50}
\definecolor{standard}{HTML}{2196F3}
\definecolor{erweiterung}{HTML}{9C27B0}

\pagestyle{fancy}
\fancyhf{}
\fancyhead[L]{\small\textcolor{grau}{Spanisch Klasse 7 -- Gesamtschule MV}}
\fancyhead[R]{\small\textcolor{grau}{LE-06a | Tier 1}}
\fancyfoot[C]{\small\thepage/\pageref{LastPage}}
\renewcommand{\headrulewidth}{0.5pt}

\newtcolorbox{kompetenzbox}{
    colback=hellgrau,
    colframe=spanisch,
    fonttitle=\bfseries,
    title=Kompetenzziele,
    arc=0pt,
    boxrule=1pt
}

\newtcolorbox{diffbox}{
    colback=white,
    colframe=grau,
    fonttitle=\bfseries,
    title=Differenzierung (nur Lehrkraft),
    arc=0pt,
    boxrule=0.5pt
}

\newtcolorbox{reflexbox}{
    colback=white,
    colframe=spanisch!50,
    fonttitle=\bfseries,
    title=Reflexion (nach der Stunde),
    arc=0pt,
    boxrule=0.5pt
}

\newcommand{\basis}{\textcolor{basis}{\textbf{[B]}}}
\newcommand{\standard}{\textcolor{standard}{\textbf{[S]}}}
\newcommand{\erweiterung}{\textcolor{erweiterung}{\textbf{[E]}}}

\begin{document}

% ============================================================================
% TITELSEITE
% ============================================================================
\begin{center}
{\Large\textbf{\textcolor{spanisch}{Langentwurf (Tier 1)}}}\\[0.5em]
{\LARGE\textbf{Los lugares -- Orte in der Stadt}}\\[0.3em]
{\large Sequenz 6a -- Mi ciudad, Std. 1--2}
\end{center}

\vspace{1em}

% ============================================================================
% KOPFBEREICH
% ============================================================================
\section*{Unterrichtsdaten}

\begin{tabularx}{\textwidth}{|l|X|l|X|}
\hline
\textbf{Fach} & Spanisch (2. FS) & \textbf{Datum} & \\
\hline
\textbf{Klasse} & 7 (Gesamtschule) & \textbf{Zeit} & 2 $\times$ 45 Min. \\
\hline
\textbf{Thema} & Orte in der Stadt & \textbf{Raum} & \\
\hline
\textbf{Lehrwerk} & Qué pasa 1, Unidad 6 & \textbf{Profil} & A (heterogen) \\
\hline
\textbf{Sequenz} & 6 von 8 (Mi ciudad) & \textbf{SuS} & ca. 25 \\
\hline
\end{tabularx}

\vspace{1em}

% ============================================================================
% LERNVORAUSSETZUNGEN
% ============================================================================
\section*{Lernvoraussetzungen}

\begin{tabularx}{\textwidth}{|l|X|}
\hline
\textbf{Sprachlich} & Grundwortschatz (Familie, Schule, Freizeit); Verben \textit{ser}, \textit{estar}, \textit{tener}, \textit{ir}; bestimmte/unbestimmte Artikel \\
\hline
\textbf{Methodisch} & Partnerarbeit, Gruppenpuzzle, Vokabellernen bekannt \\
\hline
\textbf{Inhaltlich} & Ortsangaben mit \textit{estar} + Ortsadverbien (aquí, allí) bekannt \\
\hline
\textbf{Heterogenität} & BR 20\%, MR 50\%, GY 30\% -- große Spannbreite \\
\hline
\end{tabularx}

\vspace{1em}

% ============================================================================
% KOMPETENZZIELE
% ============================================================================
\begin{kompetenzbox}

\textbf{Funktionale kommunikative Kompetenz:}
\begin{itemize}[nosep]
    \item \textbf{Hörverstehen:} Ortsangaben in Dialogen verstehen
    \item \textbf{Sprechen:} Orte in der Stadt benennen und beschreiben
    \item \textbf{Leseverstehen:} Stadtplanbeschreibungen verstehen
\end{itemize}

\textbf{Sprachliche Mittel:}
\begin{itemize}[nosep]
    \item Wortschatz: el parque, la plaza, el cine, el museo, la biblioteca, el supermercado, la farmacia, el hospital, la estación, el centro comercial
    \item Artikel: Genus bei Ortsbezeichnungen (maskulin/feminin)
    \item Strukturen: Vorbereitung auf \textit{ir a} + Ort
\end{itemize}

\textbf{Interkulturelle kommunikative Kompetenz:}
\begin{itemize}[nosep]
    \item Typische spanische Stadtbilder kennenlernen (Plaza Mayor, Mercado)
    \item Unterschiede Stadt vs. Dorf in Spanien/Deutschland
\end{itemize}

\textbf{Text- und Medienkompetenz:}
\begin{itemize}[nosep]
    \item Stadtpläne lesen und verstehen
    \item Bildbeschreibungen anfertigen
\end{itemize}

\textbf{Sprachlernkompetenz:}
\begin{itemize}[nosep]
    \item Internationalismen erkennen (hospital, museo, farmacia)
    \item Eselsbrücken für Genus entwickeln
\end{itemize}

\end{kompetenzbox}

\vspace{1em}

% ============================================================================
% STUNDENZIELE
% ============================================================================
\section*{Stundenziele}

\textbf{Stunde 1 -- Orte kennenlernen:}
\begin{enumerate}[nosep]
    \item Die SuS können 10 Orte in der Stadt auf Spanisch benennen.
    \item Die SuS können den bestimmten Artikel korrekt zuordnen.
    \item Die SuS erkennen Internationalismen und erschließen deren Bedeutung.
\end{enumerate}

\textbf{Stunde 2 -- Orte anwenden:}
\begin{enumerate}[nosep]
    \item Die SuS können Orte mit Bildern/Funktionen verknüpfen.
    \item Die SuS können einfache Sätze über Orte formulieren.
    \item Die SuS können einen Stadtplan beschreiben.
\end{enumerate}

\newpage

% ============================================================================
% VERLAUFSPLANUNG STUNDE 1
% ============================================================================
\section*{Verlaufsplanung Stunde 1: Orte kennenlernen}

\begin{tabularx}{\textwidth}{|c|c|X|c|c|l|}
\hline
\textbf{Phase} & \textbf{Zeit} & \textbf{Unterrichtsgeschehen} & \textbf{SF} & \textbf{Ziel} & \textbf{Material} \\
\hline
\multicolumn{6}{|l|}{\textbf{EINSTIEG}} \\
\hline
Begrüßung & 3' &
\begin{minipage}[t]{\linewidth}
\vspace{2pt}
L: ``¡Buenos días! Hoy hablamos de la ciudad.''\\
Zeigt Stadtbild von Madrid/Barcelona
\vspace{2pt}
\end{minipage} & PL & Neugier & Bild \\
\hline
Impuls & 5' &
\begin{minipage}[t]{\linewidth}
\vspace{2pt}
L: ``¿Qué hay en una ciudad?''\\
SuS sammeln deutsch, L führt spanisch ein\\
Bekannte Wörter aktivieren
\vspace{2pt}
\end{minipage} & PL & Vorwissen & Tafel \\
\hline
\multicolumn{6}{|l|}{\textbf{ERARBEITUNG}} \\
\hline
Input & 10' &
\begin{minipage}[t]{\linewidth}
\vspace{2pt}
L präsentiert 10 Orte mit Bildkarten:\\
el parque, la plaza, el cine, el museo,\\
la biblioteca, el supermercado, la farmacia,\\
el hospital, la estación, el centro comercial\\
SuS sprechen nach (Chor)
\vspace{2pt}
\end{minipage} & PL & Vokab. & Bildkarten \\
\hline
Sortieren & 10' &
\begin{minipage}[t]{\linewidth}
\vspace{2pt}
\textbf{Partnerarbeit:} Sortiere nach Artikel\\
Tabelle: el (maskulin) | la (feminin)\\
SuS diskutieren, finden Muster
\vspace{2pt}
\end{minipage} & PA & Genus & AB \\
\hline
\multicolumn{6}{|l|}{\textbf{SICHERUNG}} \\
\hline
Austausch & 7' &
\begin{minipage}[t]{\linewidth}
\vspace{2pt}
Ergebnisse zusammentragen:\\
Regel: -o oft maskulin, -a oft feminin\\
Ausnahmen: el día, la mano (bekannt)
\vspace{2pt}
\end{minipage} & PL & Sichern & Tafel \\
\hline
Merkkasten & 7' &
\begin{minipage}[t]{\linewidth}
\vspace{2pt}
SuS übertragen Vokabeln in AB-06a Merkkasten\\
Mit korrektem Artikel
\vspace{2pt}
\end{minipage} & EA & Sichern & AB \\
\hline
Ausblick & 3' &
\begin{minipage}[t]{\linewidth}
\vspace{2pt}
L: ``Nächste Stunde: ¿Adónde vas?''\\
Verabschiedung: ``¡Hasta luego!''
\vspace{2pt}
\end{minipage} & PL & Motivat. & -- \\
\hline
\end{tabularx}

\vspace{1em}

% ============================================================================
% VERLAUFSPLANUNG STUNDE 2
% ============================================================================
\section*{Verlaufsplanung Stunde 2: Orte anwenden}

\begin{tabularx}{\textwidth}{|c|c|X|c|c|l|}
\hline
\textbf{Phase} & \textbf{Zeit} & \textbf{Unterrichtsgeschehen} & \textbf{SF} & \textbf{Ziel} & \textbf{Material} \\
\hline
\multicolumn{6}{|l|}{\textbf{EINSTIEG}} \\
\hline
Begrüßung & 3' &
\begin{minipage}[t]{\linewidth}
\vspace{2pt}
L: ``¡Hola! ¿Qué lugares conocéis?''\\
Schnelle Wiederholung (Blitzabfrage)
\vspace{2pt}
\end{minipage} & PL & Aktivier. & -- \\
\hline
Wiederh. & 7' &
\begin{minipage}[t]{\linewidth}
\vspace{2pt}
Bildkarten-Quiz: L zeigt Bild, SuS nennen Ort\\
Mit Artikel! (``¡El parque!'')
\vspace{2pt}
\end{minipage} & PL & Festigen & Bildkarten \\
\hline
\multicolumn{6}{|l|}{\textbf{ERARBEITUNG}} \\
\hline
Funktionen & 10' &
\begin{minipage}[t]{\linewidth}
\vspace{2pt}
L: ``¿Qué hacemos en el cine?''\\
SuS ordnen Aktivitäten zu Orten:\\
cine = ver películas, biblioteca = leer libros\\
supermercado = comprar comida
\vspace{2pt}
\end{minipage} & PA & Zuordn. & AB \\
\hline
Stadtplan & 12' &
\begin{minipage}[t]{\linewidth}
\vspace{2pt}
Stadtplan von ``Villa Ejemplo'':\\
SuS lokalisieren Orte\\
``¿Dónde está el museo?''\\
``El museo está aquí.'' (zeigen)
\vspace{2pt}
\end{minipage} & PA/GA & Anwenden & Stadtplan \\
\hline
\multicolumn{6}{|l|}{\textbf{ÜBUNG}} \\
\hline
AB-Arbeit & 10' &
\begin{minipage}[t]{\linewidth}
\vspace{2pt}
AB-06a Aufgaben 2--4:\\
Zuordnung, Lückentext, kurze Sätze
\vspace{2pt}
\end{minipage} & EA & Üben & AB \\
\hline
\multicolumn{6}{|l|}{\textbf{SICHERUNG}} \\
\hline
Besprechung & 5' &
\begin{minipage}[t]{\linewidth}
\vspace{2pt}
Ergebnisse vergleichen\\
Häufige Fehler besprechen
\vspace{2pt}
\end{minipage} & PL & Korrigier. & -- \\
\hline
Abschluss & 3' &
\begin{minipage}[t]{\linewidth}
\vspace{2pt}
L: ``La próxima vez: ir + a + lugar''\\
HA: Vokabeln lernen + LP-06a\\
``¡Adiós!''
\vspace{2pt}
\end{minipage} & PL & Abschl. & -- \\
\hline
\end{tabularx}

\newpage

% ============================================================================
% DIFFERENZIERUNG
% ============================================================================
\section*{Differenzierung}

\begin{diffbox}

\textbf{Stunde 1 -- Orte kennenlernen:}

\begin{tabularx}{\textwidth}{|c|X|X|}
\hline
\textbf{Niveau} & \textbf{Maßnahme} & \textbf{Umsetzung} \\
\hline
\basis & Nur 6 Kernvokabeln (parque, cine, hospital, supermercado, farmacia, estación) & Reduzierte Liste \\
\hline
\standard & Alle 10 Vokabeln & Vollständige Aufgabe \\
\hline
\erweiterung & + 3 weitere Orte: el teatro, el ayuntamiento, la iglesia & Zusatzliste \\
\hline
\end{tabularx}

\vspace{1em}

\textbf{Stunde 2 -- Orte anwenden:}

\begin{tabularx}{\textwidth}{|c|X|X|}
\hline
\textbf{Niveau} & \textbf{Maßnahme} & \textbf{Umsetzung} \\
\hline
\basis & Zuordnung mit Bildunterstützung & Bilder auf AB \\
\hline
\standard & Zuordnung + einfache Sätze & Standard-AB \\
\hline
\erweiterung & Stadtbeschreibung (4--5 Sätze) & AB Aufgabe 4 erweitert \\
\hline
\end{tabularx}

\end{diffbox}

\vspace{1em}

% ============================================================================
% ERWARTETE SCHWIERIGKEITEN
% ============================================================================
\section*{Erwartete Schwierigkeiten}

\begin{tabularx}{\textwidth}{|l|X|X|}
\hline
\textbf{Bereich} & \textbf{Problem} & \textbf{Reaktion} \\
\hline
Genus & ``la cine'' statt ``el cine'' & Merkhilfe: kurzes Wort = oft maskulin \\
\hline
Aussprache & ``biblioteca'' [bi-bli-o-te-ka] & Silbenweise sprechen, langsam steigern \\
\hline
Orthographie & ``farmacia'' mit ``ph'' oder ``z'' & Vergleich: pharmacy (EN) vs. farmacia (ES) \\
\hline
Bedeutung & ``estación'' vs. ``estación de tren'' & Kontextualisierung: immer mit Bild \\
\hline
Länge & ``centro comercial'' als zwei Wörter & Zusammen lernen, ``Einkaufszentrum'' \\
\hline
\end{tabularx}

\vspace{1em}

% ============================================================================
% TAFELANSCHRIEB
% ============================================================================
\section*{Tafelanschrieb}

\begin{center}
\fbox{\parbox{0.9\textwidth}{
\begin{center}
\textbf{\large Los lugares de la ciudad}\\
\textbf{Orte in der Stadt}
\end{center}

\vspace{0.5em}
\begin{tabular}{l|l}
\textbf{el (maskulin)} & \textbf{la (feminin)} \\
\hline
el parque (der Park) & la plaza (der Platz) \\
el cine (das Kino) & la biblioteca (die Bibliothek) \\
el museo (das Museum) & la farmacia (die Apotheke) \\
el supermercado (der Supermarkt) & la estación (der Bahnhof) \\
el hospital (das Krankenhaus) & \\
el centro comercial (das Einkaufszentrum) & \\
\end{tabular}

\vspace{0.5em}
\textit{Tipp: Wörter auf -o sind oft maskulin, auf -a oft feminin!}
}}
\end{center}

\vspace{1em}

% ============================================================================
% MATERIALCHECKLISTE
% ============================================================================
\section*{Materialcheckliste}

\begin{tabularx}{\textwidth}{|l|X|c|c|}
\hline
\textbf{Material} & \textbf{Beschreibung} & \textbf{Std.} & \textbf{Anzahl} \\
\hline
Qué pasa 1 & Schülerbuch, Unidad 6 & 1--2 & Klassensatz \\
\hline
AB-06a & Arbeitsblatt Orte & 1--2 & Klassensatz \\
\hline
Bildkarten & 10 Orte als Karten (A4) & 1--2 & 1 Satz \\
\hline
Stadtplan & ``Villa Ejemplo'' (Fantasiestadt) & 2 & Klassensatz \\
\hline
Stadtbild & Foto Madrid/Barcelona & 1 & 1 \\
\hline
\end{tabularx}

\vspace{1em}

% ============================================================================
% HAUSAUFGABE
% ============================================================================
\section*{Hausaufgabe}

\textbf{Nach Stunde 1:} AB-06a Merkkasten vervollständigen

\textbf{Nach Stunde 2:} Vokabeln lernen (10 Orte mit Artikel); Optional: LP-06a bearbeiten

\vspace{1em}

% ============================================================================
% REFLEXION
% ============================================================================
\begin{reflexbox}

\textbf{Nach der Durchführung ausfüllen:}

\begin{enumerate}
    \item Wurde das Stundenziel erreicht?
    \begin{itemize}[nosep]
        \item[$\square$] vollständig
        \item[$\square$] größtenteils
        \item[$\square$] teilweise
        \item[$\square$] nicht erreicht
    \end{itemize}

    \item Wie war das Zeitmanagement?
    \begin{itemize}[nosep]
        \item[$\square$] passte genau
        \item[$\square$] zu wenig Zeit für: \rule{5cm}{0.4pt}
        \item[$\square$] zu viel Zeit für: \rule{5cm}{0.4pt}
    \end{itemize}

    \item Welche Schwierigkeiten traten auf?\\[2em]

    \item Was würde ich beim nächsten Mal anders machen?\\[2em]

    \item Differenzierung: War die Abstufung angemessen?
    \begin{itemize}[nosep]
        \item[$\square$] ja
        \item[$\square$] Basis zu schwer/leicht
        \item[$\square$] Erweiterung zu schwer/leicht
    \end{itemize}
\end{enumerate}

\end{reflexbox}

\vspace{1em}

% ============================================================================
% LERNSTANDSERHEBUNG
% ============================================================================
\section*{Lernstandserhebung}

\textbf{Beobachtungskriterien während des Unterrichts:}

\begin{tabularx}{\textwidth}{|X|c|c|c|}
\hline
\textbf{Kriterium} & \textbf{++} & \textbf{+} & \textbf{--} \\
\hline
SuS können alle 10 Orte mit Artikel benennen & & & \\
\hline
SuS unterscheiden el/la korrekt & & & \\
\hline
SuS sprechen die Vokabeln verständlich aus & & & \\
\hline
SuS beteiligen sich aktiv am Stadtplan-Quiz & & & \\
\hline
\end{tabularx}

\vspace{0.5em}

\textbf{Formative Überprüfung:} LP-06a (optional) oder AB-06a Aufgaben

\vspace{1em}

% ============================================================================
% ANSCHLUSS
% ============================================================================
\section*{Anschluss: Stunde 3--4 (06b)}

\textbf{Thema:} \textit{Ir a} + Ortsangaben + \textit{estar en}

\textbf{Voraussetzung:} Orte sicher beherrscht (mit Artikel)

\textbf{Neuer Inhalt:}
\begin{itemize}[nosep]
    \item Voy al cine. (ir a + el = al)
    \item Estoy en el parque.
    \item ¿Adónde vas? -- Voy a la biblioteca.
\end{itemize}

% ============================================================================
% DIDAKTIK-SCORE
% ============================================================================
\newpage
\section*{Didaktik-Score}

\begin{center}
\begin{tabular}{|c|l|c|c|p{5cm}|}
\hline
\textbf{\#} & \textbf{Dimension} & \textbf{Gew.} & \textbf{Score} & \textbf{Notiz} \\
\hline
1 & Differenzierung & 15\% & 9/10 & 3 Niveaus, qualitativ verschieden \\
\hline
2 & Taxonomie (AFB) & 10\% & 9/10 & I: benennen, II: zuordnen/anwenden, III: beschreiben \\
\hline
3 & Lernziel-Alignment & 10\% & 10/10 & Rahmenplan MV Kl. 7 konform \\
\hline
4 & Aktivierung & 10\% & 9/10 & Partnerarbeit, Bildkarten-Quiz, Stadtplan \\
\hline
5 & Scaffolding & 10\% & 9/10 & Bildunterstützung, Genus-Regel, Artikel \\
\hline
6 & Feedback-Qualität & 10\% & 8/10 & LP mit automatischem Feedback \\
\hline
7 & Sprachliche Klarheit & 10\% & 10/10 & Altersgerecht, kurze Anweisungen \\
\hline
8 & Motivation \& Relevanz & 10\% & 9/10 & Alltagsbezug: Wo gehst du hin? \\
\hline
9 & Inklusion (UDL) & 10\% & 9/10 & Bilder + Text + Audio (LP) \\
\hline
10 & Praktikabilität & 5\% & 10/10 & 90 Min. realistisch \\
\hline
\multicolumn{3}{|r|}{\textbf{GESAMT:}} & \textbf{9.2/10} & \textcolor{basis}{\textbf{FREIGABE}} \\
\hline
\end{tabular}
\end{center}

\end{document}
